
\documentclass[10pt,a4paper]{article}

\usepackage[utf8]{inputenc}
\usepackage[T1]{fontenc}
\usepackage[french]{babel}
\usepackage{amsmath,amssymb}
\usepackage{geometry}
\usepackage{enumitem}
\usepackage{multicol}
\usepackage{tikz}
\usepackage{tasks}
\usepackage{pgfplots}

\geometry{
    top=1.5cm,
    bottom=1.5cm,
    left=1.7cm,
    right=1.7cm
}

\setlength{\parindent}{0pt}

\newcommand{\version}[1]{
    \begin{center}
        \textbf{\large INTERROGATION -- Fonctions et proportions}\\
        \textbf{Version #1}
    \end{center}

    \vspace{-0.1cm}

    \textbf{Nom :} \rule{5cm}{0.4pt}
    \hfill
    \textbf{Prénom :} \rule{4cm}{0.4pt}

    \vspace{0.2cm}
}

\begin{document}

% =========================================================
% VERSION A
% =========================================================

\version{A}

\textbf{Exercice 1 -- DC \hfill /2}

Soit \(g\) la fonction définie sur \(\mathbb{R}\) par \(g(x)=x^2-3x+2\). Calculer \(g(-2)\).

\vspace{0.1cm}

\textbf{Exercice 2 -- Lecture graphique \hfill /4}

On considère une fonction \(f\) dont la représentation graphique est donnée ci-dessous.
\begin{multicols}{2}
    \begin{enumerate}
    \item Quel est l'ensemble de définition de \(f\) ?
    \item Donner l'image de -2 par \(f\).
    \item Donner les éventuels antécédents de -1
    \item Résoudre \(f(x)=0\).
    \end{enumerate}

    \begin{tikzpicture}[scale=0.8]
    \begin{axis}[
        axis lines=middle,
        xmin=-2.5, xmax=2.9,
        ymin=-6, ymax=2,
        domain=-2:2,
        samples=300,
        grid=major,
        xtick={-2,-1,0,1,2,3},
        ytick={-5,-4,-3,-2,-1,0,1,2,3,4},
        xlabel={\(x\)},
        ylabel={\(y\)},
        xlabel style={
            at={(axis cs:2.9,0)},
            anchor=west
        },
        ylabel style={
            at={(axis cs:0,2)},
            anchor=south
        },
        width=9cm,
        height=8cm
    ]

    \addplot[
        blue,
        very thick
    ] {-2/3*(x-1)*(x+2.5)*(x-0)+1}
    node[pos=0.92, above right] {\(\mathcal{C}_f\)};

    \end{axis}
    \end{tikzpicture}

\end{multicols}

\textbf{Exercice 3 -- fonction affine \hfill /4}

Soit \(h\) la fonction affine définie par \(h(x)=2x-3\). Déterminer son tableau de signe.

\vfill

% =========================================================
% VERSION B
% =========================================================

\version{B}

\textbf{Exercice 1 -- DC \hfill /2}

Soit \(g\) la fonction définie sur \(\mathbb{R}\) par \(g(x)=2(4-x)(x+1)\). Calculer \(g(-2)\).

\vspace{0.1cm}

\textbf{Exercice 2 -- Lecture graphique \hfill /4}

On considère une fonction \(f\) dont la représentation graphique est donnée ci-dessous.
\begin{multicols}{2}
    \begin{enumerate}
    \item Quel est l'ensemble de définition de \(f\) ?
    \item Donner l'image de 2 par \(f\).
    \item Donner les éventuels antécédents de 1
    \item Résoudre \(f(x)=0\).
    \end{enumerate}

    \begin{tikzpicture}[scale=0.8]
    \begin{axis}[
        axis lines=middle,
        xmin=-2.5, xmax=2.9,
        ymin=-4, ymax=2,
        domain=-2:2,
        samples=300,
        grid=major,
        xtick={-2,-1,0,1,2,3},
        ytick={-5,-4,-3,-2,-1,0,1,2,3,4},
        xlabel={\(x\)},
        ylabel={\(y\)},
        width=9cm,
        height=8cm,
        xlabel style={
            at={(axis cs:2.9,0)},
            anchor=west
        },
        ylabel style={
            at={(axis cs:0,2)},
            anchor=south
        },
    ]

    \addplot[
        blue,
        very thick
    ] {-2/3*(x/1.2-1)*(x/1.2+2.5)*(x/1.2-0)+0.5}
    node[pos=0.92, above right] {\(\mathcal{C}_f\)};

    \end{axis}
    \end{tikzpicture}

\end{multicols}

\textbf{Exercice 3 -- fonction affine \hfill /4}

Soit \(h\) la fonction affine définie par \(h(x)=-4x-3\). Déterminer son tableau de signe.

\newpage

% =========================================================
% VERSION C
% =========================================================

\version{C}

\textbf{Exercice 1 -- DC \hfill /2}

Soit \(g\) la fonction définie sur \(\mathbb{R}\) par \(g(x)=x^2+3x+2\). Calculer \(g(-5)\).

\vspace{0.1cm}

\textbf{Exercice 2 -- Lecture graphique \hfill /4}

On considère une fonction \(f\) dont la représentation graphique est donnée ci-dessous.
\begin{multicols}{2}
    \begin{enumerate}
    \item Quel est l'ensemble de définition de \(f\) ?
    \item Donner l'image de 1 par \(f\).
    \item Donner les éventuels antécédents de -1
    \item Résoudre \(f(x)=0\).
    \end{enumerate}

    \begin{tikzpicture}[scale=0.8]
    \begin{axis}[
        axis lines=middle,
        xmin=-3.5, xmax=1.9,
        ymin=-2, ymax=5.5,
        domain=-3:1,
        samples=300,
        grid=major,
        xtick={-3,-2,-1,0,1,2},
        ytick={-5,-4,-3,-2,-1,0,1,2,3,4,5},
        xlabel={\(x\)},
        ylabel={\(y\)},
        xlabel style={
            at={(axis cs:1.9,0)},
            anchor=west
        },
        ylabel style={
            at={(axis cs:0,5.5)},
            anchor=south
        },
        width=9cm,
        height=8cm
    ]

    \addplot[
        blue,
        very thick
    ] {-2/3*(x-1)*(x+2.5)*(x-0)+1}
    node[pos=0.92, above right] {\(\mathcal{C}_f\)};

    \end{axis}
    \end{tikzpicture}

\end{multicols}

\textbf{Exercice 3 -- fonction affine \hfill /4}

Soit \(h\) la fonction affine définie par \(h(x)=2x+3\). Déterminer son tableau de signe.

\vfill

% =========================================================
% VERSION D
% =========================================================

\version{D}
\textbf{Exercice 1 -- DC \hfill /2}

Soit \(g\) la fonction définie sur \(\mathbb{R}\) par \(g(x)=2(x-2)(x+2)\). Calculer \(g(-3)\).

\vspace{0.1cm}

\textbf{Exercice 2 -- Lecture graphique \hfill /4}

On considère une fonction \(f\) dont la représentation graphique est donnée ci-dessous.
\begin{multicols}{2}
    \begin{enumerate}
    \item Quel est l'ensemble de définition de \(f\) ?
    \item Donner l'image de 3 par \(f\).
    \item Donner les éventuels antécédents de -1
    \item Résoudre \(f(x)=0\).
    \end{enumerate}

    \begin{tikzpicture}[scale=0.8]
    \begin{axis}[
        axis lines=middle,
        xmin=-0.5, xmax=4.9,
        ymin=-2, ymax=4,
        domain=0:4,
        samples=300,
        grid=major,
        xtick={-2,-1,0,1,2,3,4},
        ytick={-5,-4,-3,-2,-1,0,1,2,3,4},
        xlabel={\(x\)},
        ylabel={\(y\)},
        xlabel style={
            at={(axis cs:4.9,0)},
            anchor=west
        },
        ylabel style={
            at={(axis cs:0,4)},
            anchor=south
        },
        width=9cm,
        height=8cm
    ]

    \addplot[
        blue,
        very thick
    ] {1/30*(14*x^3-63*x^2+43*x)+1}
    node[pos=0.92, above right] {\(\mathcal{C}_f\)};

    \end{axis}
    \end{tikzpicture}

\end{multicols}

\textbf{Exercice 3 -- fonction affine \hfill /4}

Soit \(h\) la fonction affine définie par \(h(x)=-4x+3\). Déterminer son tableau de signe.

\end{document}

